% Рубежный контроль №1, вариант 1 — по материалам: лекции (mathmod.pdf),
% методические указания (23\_03\_2020\_...\_Моделирование\_2.doc),
% вопросы/конспект (Modelirovanie\_1\_-3.pdf).
\documentclass[12pt,a4paper]{article}
\usepackage[T2A]{fontenc}
\usepackage[utf8]{inputenc}
\usepackage[russian]{babel}
\usepackage{amsmath,amssymb}
\usepackage{geometry}
\geometry{margin=2.2cm}
\usepackage{enumitem}
\setlist{nosep,leftmargin=*}

\title{Рубежный контроль №1 (вариант 1)\\Ответы}
\author{}
\date{}

\begin{document}
\maketitle

\begin{enumerate}[label=\textbf{\arabic*.}]

\item \textbf{Аналитическое моделирование (определение).}

Аналитическая модель в смысле курса — это \emph{формализация явления или процесса реального мира, учитывающая физико-химические воздействия, протекающие в системе}. Аналитическое моделирование — построение и использование таких моделей, когда связи между величинами задаются в аналитическом (формульном) виде и допускают теоретический анализ и расчёт на ЭВМ.

\item \textbf{Обоснование выбора метода построения цифровой модели рельефа.}

По методическим указаниям к лабораторной работе: ставится задача \emph{построить модель фрагмента поверхности по заданной матрице высот} и \emph{обосновать выбор численного метода для приближённого вычисления высоты в заданной точке}. В качестве \emph{цифровой модели поверхности} задаётся матрица высот; в качестве метода интерполяции предлагается выбрать \emph{одну из модификаций триангуляции Делоне}. Обоснование такого выбора в рамках того же источника: триангуляция Делоне даёт хорошо сбалансированную сетку (в предельном случае треугольники стремятся к равносторонним), что снижает искажения при интерполяции; для произвольной точки плоскости высота находится по узлам ограниченного числа соседних треугольников; для метода и его модификаций разработаны эффективные алгоритмы с изучаемой вычислительной сложностью (см.\ также рекомендуемую в указаниях литературу по триангуляции Делоне). Сравнение с альтернативами (например, только регулярная сетка без учёта рельефа, глобальные сплайновые поверхности) проводят по точности на контрольных точках, устойчивости и затратам на хранение и пересчёт.

\item \textbf{Имитационная модель (определение).}

В методических указаниях и конспекте \texttt{Modelirovanie\_1\_-3.pdf}: имитационная модель \emph{представляет собой численный метод вычислительного эксперимента} и \emph{позволяет имитировать поведение реальной системы на ЦВМ} в заданный или формируемый период времени с определённым набором входных данных. Для много-компонентных систем имитационное моделирование опирается на идею «чёрного ящика» и задачи идентификации: задаётся преобразование входов (и управляющих воздействий) в выходы; детализация физико-химических процессов может быть свёрнута до закономерности изменения выходных данных.

\item \textbf{Полунатурный эксперимент (определение).}

Компромисс между натурным и вычислительным экспериментом: \emph{в качестве объекта исследования используется физическая модель реального объекта}; часть характеристик измеряется экспериментально (пример из курса — число Маха в аэродинамической трубе), другие задаются или уточняются программно. Результат может потребовать дополнительных вычислительных задач для повышения точности оценки.

\item \textbf{Постановка баллистической задачи.}

В базовой постановке: материальная точка (тело) брошена под углом $\alpha$ к горизонту со скоростью $v_0$; требуется, например, определить дальность полёта, время, траекторию. \emph{Модель Галилея:} плоская Земля, $g=\mathrm{const}$, сопротивление воздуха пренебрегается; движение задаётся параметрически $x=v_0\cos\alpha\cdot t$, $y=v_0\sin\alpha\cdot t-gt^2/2$. \emph{Модель Ньютона:} учитывается сила лобового сопротивления воздуха (например, $F=-\beta v^2$), получается система ОДУ на продольную и вертикальную составляющие скорости. В курсе отмечается, что современная баллистическая постановка обобщается (линейные и угловые параметры, до шести неизвестных и более сложная модель среды).

\item \textbf{Способы организации квазипараллелизма (перечислить).}

При дискретном модельном времени используют схемы псевдопараллелизма:
\begin{enumerate}[label=\arabic*)]
\item моделирование \textbf{активностями};
\item \textbf{событиями};
\item \textbf{транзактами};
\item \textbf{процессами};
\item \textbf{агрегатами}.
\end{enumerate}

\item \textbf{Предмет машинного обучения (указать).}

В \texttt{mathmod.pdf} термин «машинное обучение» отдельно не определён; из близкого материала: построение правила (алгоритма), сопоставляющего объекты обучающей выборки меткам кластеров; различают \emph{обучение с учителем} (классификация при заданном числе интерпретируемых кластеров) и \emph{обучение без учителя} (кластеризация). В общем виде предмет МО — \emph{математические и алгоритмические методы построения зависимостей и правил принятия решений по данным}, качество которых улучшается по мере накопления примеров (обучающей информации).

\item \textbf{Функциональная зависимость (определение).}

\emph{Функциональная зависимость} — это соответствие, при котором некоторой величине $y\in Y$ ставится в соответствие \emph{единственная} величина $x\in X$ (в обсуждаемой на лекции постановке для пар $(x,y)$ одному $x$ соответствует одно $y$; иначе говорят о стохастической зависимости и переходят к регрессии, оценке тренда и погрешности).

\item \textbf{Стохастическая модель (определение).}

По классификации из методических указаний и конспекта: \emph{стохастическая модель} — это модель, \emph{входными и выходными данными которой являются случайные величины}. Если модель опирается на основную задачу теории вероятностей, говорят о \emph{вероятностной модели}. Дополнительно в курсе: стохастические (случайные) зависимости между величинами изучаются через корреляцию, регрессию и проверку статистических гипотез; имитационное моделирование с подбором параметров по случайным реализациям входов связывают со «стохастической настройкой» системы.

\item \textbf{Построение модифицированного метода моментов (линейная функциональная зависимость).}

\emph{Постановка.} Даны пары наблюдений $(x_i,y_i)$, $i=1,\ldots,n$, и предполагается, что среднее значение отклика связано с фактором линейно:
\[
\mathbb{E}[Y\mid X=x] = \theta_0 + \theta_1 x
\]
(в терминологии курса — линейная регрессия; при случайных $X,Y$ параметры связи оценивают по выборке). Классический метод моментов для распределения сопоставляет теоретические моменты с выборочными; для оценки параметров \emph{связи} двух величин используют \emph{модификацию}: составляют уравнения, эквивалентные сопоставлению моментов совместного распределения $(X,Y)$ с эмпирическими.

Удобная пара уравнений (запись через первый момент $Y$ и ковариацию $(X,Y)$ при стандартных предположениях о модели ошибки):
\[
\mathbb{E}[Y] = \theta_0 + \theta_1\,\mathbb{E}[X], \qquad
\mathrm{cov}(X,Y) = \theta_1\,\mathrm{var}(X).
\]
Заменяя теоретические величины выборочными: $\bar{x}=\frac1n\sum x_i$, $\bar{y}=\frac1n\sum y_i$,
\[
\widehat{S}_{xx}=\frac1n\sum(x_i-\bar{x})^2,\quad
\widehat{S}_{xy}=\frac1n\sum(x_i-\bar{x})(y_i-\bar{y}),
\]
получают оценки
\[
\hat\theta_1 = \frac{\widehat{S}_{xy}}{\widehat{S}_{xx}}, \qquad
\hat\theta_0 = \bar{y} - \hat\theta_1\,\bar{x}
\]
(при $\widehat{S}_{xx}>0$). Это тот же вид оценок, что даёт метод наименьших квадратов для простой линейной регрессии; в курсе подчёркивается, что модификации метода моментов и максимального правдоподобия применяют, когда величины ненаблюдаемы одновременно и прямое МНК невозможно — тогда систему моментных уравнений подбирают под доступные данные.

\item \textbf{Условия Делоне (формулировка).}

По методическим указаниям (определение триангуляции Делоне): \emph{под триангуляцией Делоне в общем случае подразумевают такой вид триангуляции, когда описанная для каждого треугольника окружность не содержит внутри себя ни одной точки набора}. Отмечается также, что такая триангуляция \emph{хорошо сбалансирована} в том смысле, что формируемые треугольники в предельном случае стремятся к \emph{равносторонним}. Эквивалентные геометрические критерии (локальный тест на общем ребре, максиминный угол) используют при доказательствах и в программных реализациях.

Для пошагового построения на ЭВМ в тех же указаниях вводят \emph{классификацию рёбер} триангуляции Делоне относительно текущей триангуляции: ребро может быть \emph{«спящим»} (ещё не обнаружено алгоритмом), \emph{«живым»} (обнаружено, известна только одна примыкающая область) и \emph{«мёртвым»} (обработано); алгоритм завершается, когда не остаётся «живых» рёбер.

\item \textbf{Математическая постановка задачи машинного обучения.}

Обобщая подход курса к регрессии и многомерному анализу: заданы пространство объектов $\mathcal{X}$, пространство ответов $\mathcal{Y}$, неизвестная зависимость $y \approx f(x)$ и выборка $(x_i,y_i)_{i=1}^n$. Задаётся класс моделей $\mathcal{F}$ (гипотез) и функция потерь $L(y,\hat{y})$. Ищут
\[
f^* \in \arg\min_{f\in\mathcal{F}} \frac1n\sum_{i=1}^n L\bigl(y_i, f(x_i)\bigr)
\]
(эмпирический риск); часто добавляют регуляризацию. Для задач без учителя (кластеризация в курсе) минимизируют критерий качества разбиения выборки на кластеры в выбранной метрике. Теоретически целью является близость к минимуму ожидаемого риска по распределению данных, оцениваемому по выборке.

\end{enumerate}

\end{document}
